\chapter{Stand der Technik und verwandte Arbeiten}
\label{ch:std_introduction}

Dieses Kapitel analysiert den aktuellen Forschungsstand zu Forward- und Deferred-Rendering-Architekturen sowie moderne
Industriepraktiken im Bereich hybrider Architekturen und selektiver Stilisierung.
Durch die Betrachtung bestehender Verfahren im Kontext des Non-Photorealistic Renderings (NPR) werden deren
konzeptionelle und leistungstechnische Limitierungen aufgezeigt.
Die Identifikation dieser Schwachstellen sowie der Mangel an empirischen Vergleichsdaten dienen der Definition der
zugrundeliegenden Forschungslücke, legitimieren die in dieser Arbeit durchgeführte holistische Systemanalyse und
ermöglichen die wissenschaftliche Einordnung des eigens entwickelten Evaluations-Frameworks.


\section{Aktueller Forschungsstand zu Forward- und Deferred-Architekturen}
\label{sec:sdt_academic_state}

% Einleitung
Deferred-Rendering-Paradigmen stellen sowohl in der akademischen Forschung als auch in der industriellen Praxis
etablierte Standards dar~\parencite[vgl.][S. 2]{mallettDeferredAdaptiveCompute2018}.
Moderne Echtzeit-Rendering-Pipelines werden maßgeblich von komplexen Beleuchtungsanforderungen getrieben, wodurch das
effiziente Handhaben von Deferred-Architekturen in den Fokus rückt.
Gleichzeitig werden im Bereich des Forward-Renderings durch diverse Optimierungsverfahren die klassischen Schwächen
dieses Ansatzes zunehmend
relativiert~\parencite[vgl.]{olssonClusteredDeferredForward2012, haradaForwardBringingDeferred2012}.
Die Wahl des Rendering-Ansatzes ist in der Praxis stark kontextabhängig und auf die spezifischen Leistungs- und
Qualitätsanforderungen des Zielbildes ausgerichtet.
Folglich existieren sowohl Deferred- als auch Forward-Implementierungen in modernen Engines selten in ihrer reinen
Form~\parencite[vgl.][S. 56]{burnsVisibilityBufferCacheFriendly2013}.
Der wissenschaftliche Diskurs fokussiert sich daher primär auf hybride Lösungsansätze und gezielte Reduktionen
architektonischer Engpässe.

% Forward Rendering
Der primäre architektonische Nachteil naiver Forward-Rendering-Ansätze liegt in der exponentiellen Leistungsdegradation
bei hoher Geometrie- und Beleuchtungskomplexität.
Die algorithmische Komplexität skaliert mit $\mathcal{O}(G \cdot L)$, wobei $G$ die Anzahl der Primitiven und $L$ die
Anzahl der Lichtquellen darstellt~\parencite[vgl.][S. 1]{Thaler2011DeferredRendering}.
Zu dieser Auswertungskomplexität addiert sich eine kombinatorische Explosion der Shader-Permutationen, sobald diverse
Materialeigenschaften und Lichttypen in der Szene verrechnet werden müssen.
Dieser Limitierung liegt die inhärente Eigenschaft von Forward-Pipelines zugrunde, bei hoher Szenen-Entropie
zwangsweise eine gewisse Menge Fragment-Overdraw zu erzeugen~\parencite[vgl.][S. 1]{fatahalianReducingShadingGPUs2010}.
Ansätze wie ein Z-Prepass zur frühzeitigen Verwerfung verdeckter Fragmente gehören in Forward-Pipelines zwar zum
Standard~\parencite[vgl.][S. 57]{burnsVisibilityBufferCacheFriendly2013}, reichen bei hoher Objekt- und Lichtdichte
jedoch nicht aus, um die Effizienz von Deferred-Verfahren zu
erreichen~\parencite[vgl.][S. 4]{Thaler2011DeferredRendering}.
Modernere Paradigmen wie \enquote{Clustered Forward} oder\enquote{Forward+} adressieren dieses Problem, indem
Lichtquellen in 2D-Kacheln oder 3D-Clustern räumlich sortiert
werden~\parencite[vgl.]{olssonClusteredDeferredForward2012, haradaForwardBringingDeferred2012}.
Dies ermöglicht eine Effizienz, die Deferred-Ansätzen nahekommt, während der primäre Vorteil der Forward-Architektur,
die hohe Flexibilität bei der Materialauswertung, erhalten bleibt.

% Deferred Rendering
Deferred-Rendering-Pipelines haben sich in Umgebungen mit hoher geometrischer Komplexität und anspruchsvollen
Beleuchtungsszenarien weitestgehend durchgesetzt, weisen jedoch ebenfalls diverse architektonische Limitierungen
auf.
Der signifikanteste Flaschenhals ist die hohe Auslastung der Speicherbandbreite (Memory Bandwidth), welche durch den
permanenten Schreib- und Lesezugriff auf den G-Buffer induziert
wird~\parencite[vgl.][S. 55]{burnsVisibilityBufferCacheFriendly2013}.
Ein moderner Lösungsansatz ist die Implementierung entkoppelter Architekturen wie des \textit{Visibility-Buffers}
(V-Buffer)~\parencite{burnsVisibilityBufferCacheFriendly2013}, welcher als Shading-Grundlage für
hochkomplexe Geometrie-Pipelines wie der \textit{Nanite}-Architektur in der Unreal Engine 5
dient~\parencite[vgl.][Folie \enquote{Visibility Buffer}]{Karis_Nanite_SIGGRAPH_Advances_2021_final}.
Im Gegensatz zum traditionellen G-Buffer speichert der V-Buffer pro Pixel lediglich minimale geometrische Metadaten
(wie Primitive- und Instanz-IDs) sowie den Tiefenwert.
Die speicherintensive Materialauswertung wird vollständig in einen nachgelagerten Pass verschoben.
Neben der Bandbreitenproblematik existieren in klassischen Deferred-Architekturen zwei weitere Nachteile: die fehlende
native Kompatibilität mit hardwarebasiertem Multisample Anti-Aliasing (MSAA) sowie die Einschränkung bei der Darstellung
transluzenter Oberflächen.
Da die Beleuchtung im Screen-Space berechnet wird, sind die für MSAA erforderlichen Geometriedaten zu diesem Zeitpunkt
verworfen.
Stattdessen müssen temporale Post-Processing-Verfahren (wie TAA) eingesetzt
werden~\parencite[vgl.][S. 5]{Thaler2011DeferredRendering, taaantialiasing}.
Die Evaluierung semitransparenter Geometrie ist im reinen Deferred-Pfad nicht trivial lösbar, da ein Standard-G-Buffer
pro Pixel nur eine diskrete Tiefenschicht vorhält.
Dieses Defizit wird in der industriellen Praxis standardmäßig über einen nachgelagerten hybriden Forward-Pass für
transparente Objekte gelöst~\parencite[vgl.][S. 56]{burnsVisibilityBufferCacheFriendly2013}.

% NPR
Hinsichtlich des Non-Photorealistic Renderings (NPR) unterscheidet die aktuelle Literatur primär zwischen Object-Space-
und Screen-Space-Methoden~\parencite[vgl.][S. 148]{magdicsPostprocessingNPREffects2013}.
Object-Space-Ansätze sind in der Regel in Forward-Pipelines verankert und berechnen Stilisierungen direkt auf der
Geometrie, was eine granulare, objektbasierte Anpassung ermöglicht.
Bei hoher Szenenkomplexität weisen diese Verfahren jedoch eine unzureichende Skalierbarkeit auf, was primär auf den
Fragment-Overdraw bei rechenintensiven Abstraktionseffekten zurückzuführen
ist~\parencite[vgl.][S. 1]{fatahalianReducingShadingGPUs2010}.
Die moderne Rendering-Forschung weicht zur Umgehung dieser Performance-Limits oft auf Screen-Space-Methoden
(Post-Processing) aus, was architektonisch stark den Paradigmen des Deferred-Renderings entspricht.
Wie~\textcite[S. 148]{magdicsPostprocessingNPREffects2013} beschreiben, ist diese Screen-Space-Verarbeitung zwar
sehr efficient, löst jedoch den direkten Bezug zum individuellen 3D-Objekt auf, wodurch etwa die selektive Stilisierung
stark erschwert wird.
Um räumlichen Kontext zurückzugewinnen, wird in der Forschung auf hybride Ansätze wie das \enquote{2.5D NPR}
zurückgegriffen~\parencite[vgl.][S. 17]{Thaler2011DeferredRendering}.
Hierbei fungiert der G-Buffer als Datenquelle für geometrische Metadaten, um komplexe Effekte geometriebewusst im
Screen-Space durchzuführen.
Klassische Deferred-Pipelines weisen in diesem Kontext jedoch eine fundamentale Schwäche bezüglich der Materialvielfalt
auf~\parencite[vgl.][S. 1]{haradaForwardBringingDeferred2012}.
Wie~\textcite[Kap. 2.6]{Thaler2011DeferredRendering} beschreibt, ist das Auswerten zahlreicher spezifischer Materialien
oder Stilisierungen in einem vereinheitlichten Screen-Space-Pass nur unter starken Performance-Einbußen durch
\textit{Dynamic Branching} auf der GPU realisierbar.
In der Konsequenz bleiben performante Stilisierungen bei reinem Deferred-Rendering zumeist auf globale
Post-Processing-Filter beschränkt~\parencite[vgl.][S. 148]{magdicsPostprocessingNPREffects2013}.

% Fazit
Zusammenfassend lässt sich festhalten, dass die aktuelle Forschung zwar diverse architektonische Lösungen für hohe
Geometrie- und Beleuchtungskomplexität sowohl in Forward- als auch in Deferred-Architekturen bietet, die spezifische
Schnittmenge mit selektivem Non-Photorealistic Rendering in Deferred-Pipelines jedoch theoretischen Limitierungen
unterliegt.
Um zu evaluieren, ob und wie diese akademischen Hürden in der Praxis umgangen werden, ist im Folgenden ein genauerer
Blick auf die angewandte Industriepraxis moderner Game-Engines erforderlich.


\section{Verfahren zur selektiven Stilisierung in der Industriepraxis}
\label{sec:sdt_industry_state}


\section{Limitierungen bestehender Ansätze und Identifikation der Forschungslücke}
\label{sec:sdt_usp}